\documentclass[11pt,a4paper]{article}
\usepackage[utf8]{inputenc}
\usepackage[vietnamese]{babel}
\usepackage{amsmath,amssymb}
\usepackage{geometry}
\geometry{a4paper, total={180mm,260mm}, left=15mm, right=15mm, top=18mm, bottom=18mm}
\usepackage{tikz}
\usetikzlibrary{shapes.geometric, arrows.meta, positioning, calc}
\usepackage{tcolorbox}
\tcbuselibrary{skins,breakable}
\usepackage{xcolor}
\usepackage{multicol}
\usepackage{array}
\usepackage{fontawesome}
\usepackage{fancyhdr}

\pagestyle{fancy}
\fancyhf{}
\renewcommand{\headrulewidth}{0.5pt}
\renewcommand{\footrulewidth}{0.5pt}
\fancyhead[L]{\small\bfseries\color{teal}\faBook\ SỔ TAY KIẾN THỨC TOÁN 3}
\fancyhead[R]{\small\bfseries\color{teal}Lớp Toán Thầy Lĩnh -- loptoanthaylinh.vn}
\fancyfoot[C]{\small\color{gray}-- Sổ tay kiến thức: Bài 58. Luyện tập chung: Phép tính trong phạm vi 10 000 --}

\definecolor{primary}{HTML}{005F73}
\definecolor{secondary}{HTML}{0A9396}
\definecolor{accent}{HTML}{EE9B00}
\definecolor{danger}{HTML}{AE2012}
\definecolor{boxbg}{HTML}{F4F6F8}

\begin{document}

% --- BANNER TIÊU ĐỀ ---
\begin{tcolorbox}[
  enhanced,
  colback=white,
  colframe=primary,
  arc=3mm,
  boxrule=1.2pt,
  left=4mm, right=4mm, top=2.5mm, bottom=2.5mm,
  drop shadow=black!10!white
]
  \centering
  {\Large\bfseries\color{primary}\faBookmark\ BÀI 58: LUYỆN TẬP CHUNG: PHÉP TÍNH TRONG PHẠM VI 10 000}\\[1.5mm]
  {\small\itshape\color{black!75}Đại tổng kết 4 phép tính, tìm thành phần chưa biết, tính giá trị biểu thức và giải toán 2 bước trên 1 trang A4}
\end{tcolorbox}

\vspace{-2mm}


\noindent
\begin{minipage}[t]{0.49\textwidth}\vspace{0pt}
\begin{tcolorbox}[
  enhanced,
  colback=blue!4!white,
  colframe=blue!75!black,
  arc=2mm, boxrule=0.8pt,
  left=2mm, right=2mm, top=1.8mm, bottom=1.8mm,
  title={\footnotesize\bfseries\color{white}\faCalculator\ 1. Đại tổng kết 4 phép tính}
]
{\footnotesize
\begin{itemize}\setlength{\itemsep}{1.5pt}\setlength{\leftskip}{-2.5mm}
  \item \textbf{Cộng và Trừ:} Đặt tính thẳng cột, thực hiện từ phải sang trái. Nhớ và mượn chính xác.
  \item \textbf{Nhân:} Nhân thừa số 1 chữ số lần lượt từ phải sang trái, nhân trước cộng nhớ sau.
  \item \textbf{Chia:} Chia từ trái sang phải, nhớ quy tắc ghi chữ số $0$ vào thương khi hạ số bé hơn số chia.
  \item Thứ tự biểu thức: Ngoặc trước $\to$ Nhân chia $\to$ Cộng trừ.
\end{itemize}
}
\end{tcolorbox}
\end{minipage}
\hfill
\begin{minipage}[t]{0.49\textwidth}\vspace{0pt}
\begin{tcolorbox}[
  enhanced,
  colback=teal!4!white,
  colframe=teal!75!black,
  arc=2mm, boxrule=0.8pt,
  left=2mm, right=2mm, top=1.8mm, bottom=1.8mm,
  title={\footnotesize\bfseries\color{white}\faSliders\ 2. Tìm thành phần chưa biết}
]
{\footnotesize
\begin{itemize}\setlength{\itemsep}{1.5pt}\setlength{\leftskip}{-2.5mm}
  \item $x + a = b \implies x = b - a$ (Tìm số hạng).
  \item $x - a = b \implies x = b + a$ (Tìm số bị trừ).
  \item $a - x = b \implies x = a - b$ (Tìm số trừ).
  \item $x \times a = b \implies x = b : a$ (Tìm thừa số).
  \item $x : a = b \implies x = b \times a$ (Tìm số bị chia).
  \item $a : x = b \implies x = a : b$ (Tìm số chia).
\end{itemize}
}
\end{tcolorbox}
\end{minipage}

\vspace{-2.5mm}
% --- KHỐI 3: SƠ ĐỒ TƯ DUY TIKZ ---
\noindent
\begin{tcolorbox}[
  enhanced,
  colback=orange!4!white,
  colframe=orange!80!black,
  arc=2mm, boxrule=0.8pt,
  left=2mm, right=2mm, top=1.5mm, bottom=1.5mm,
  title={\footnotesize\bfseries\color{white}\faSitemap\ SƠ ĐỒ TƯ DUY: HỆ THỐNG 4 PHÉP TOÁN CỐT LÕI PHẠM VI 10 000}
]
\centering
\begin{tikzpicture}[
  box/.style={draw, rounded corners=2pt, align=center, font=\scriptsize\bfseries, inner sep=2.5pt, minimum height=0.75cm},
  arr/.style={-{Stealth[length=2mm]}, thick}
]
  \node[box, fill=purple!15, draw=purple!80!black, text width=2.6cm] (k1) at (0, 0) {PHÉP CỘNG\\Đặt tính thẳng cột\\Nhớ sang hàng trái};
  \node[box, fill=teal!15, draw=teal!80!black, text width=2.6cm] (k2) at (3.5, 0) {PHÉP TRỪ\\Mượn 1 chục\\Trả nhớ sang số trừ};
  \node[box, fill=blue!12, draw=blue!80!black, text width=2.6cm] (k3) at (7.0, 0) {PHÉP NHÂN\\Nhân phải qua trái\\Nhân trước cộng nhớ};
  \node[box, fill=orange!15, draw=orange!80!black, text width=2.6cm] (k4) at (10.5, 0) {PHÉP CHIA\\Chia trái qua phải\\Xử lý chữ số 0};

  \draw[arr, purple!80!black] (k1) -- (k2);
  \draw[arr, teal!80!black] (k2) -- (k3);
  \draw[arr, blue!80!black] (k3) -- (k4);
\end{tikzpicture}
\end{tcolorbox}

\vspace{-2.5mm}
% --- KHỐI 4: VÍ DỤ & MẸO TRÁNH BẪY ---
\noindent
\begin{minipage}[t]{0.49\textwidth}\vspace{0pt}
\begin{tcolorbox}[
  enhanced,
  colback=green!4!white,
  colframe=green!60!black,
  arc=2mm, boxrule=0.8pt,
  left=2mm, right=2mm, top=1.8mm, bottom=1.8mm,
  title={\footnotesize\bfseries\color{white}\faCheckCircle\ Ví dụ thực hành chuẩn mực}
]
{\footnotesize
\textbf{1. Tìm $x$:}\\
$x \times 3 = 4560 \implies x = 4560 : 3 = \mathbf{1520}$.\\
$8400 - x = 2150 \implies x = 8400 - 2150 = \mathbf{6250}$.\\
\textbf{2. Tính giá trị biểu thức:}\\
$A = 1200 \times 4 - 2500 = 4800 - 2500 = \mathbf{2300}$.\\
$B = (3500 + 1500) : 5 = 5000 : 5 = \mathbf{1000}$.
}
\end{tcolorbox}
\end{minipage}
\hfill
\begin{minipage}[t]{0.49\textwidth}\vspace{0pt}
\begin{tcolorbox}[
  enhanced,
  colback=red!4!white,
  colframe=red!70!black,
  arc=2mm, boxrule=0.8pt,
  left=2mm, right=2mm, top=1.8mm, bottom=1.8mm,
  title={\footnotesize\bfseries\color{white}\faExclamationTriangle\ Mẹo tránh bẫy trắc nghiệm}
]
{\footnotesize
\begin{itemize}\setlength{\itemsep}{1.5pt}\setlength{\leftskip}{-2.5mm}
  \item \textbf{Bẫy tìm số trừ và số bị trừ:}\\
  -- Tìm số trừ: lấy số bị trừ \textbf{trừ đi} hiệu ($x = a - b$).\\
  -- Tìm số bị trừ: lấy hiệu \textbf{cộng với} số trừ ($x = b + a$).
  \item Đừng nhầm lẫn giữa hai quy tắc này khi làm bài tìm $x$!
\end{itemize}
}
\end{tcolorbox}
\end{minipage}


\end{document}
